\documentclass{beamer}

\mode<presentation>
{
  \usetheme{Singapore}      % or try Darmstadt, Madrid, Warsaw, ...
  \usecolortheme{crane} % or try albatross, beaver, crane, ...
  \usefonttheme{serif}  % or try serif, structurebold, ...
  \setbeamertemplate{navigation symbols}{}
  \setbeamertemplate{caption}[numbered]
} 

\usepackage[english]{babel}
\usepackage[autostyle=true]{csquotes}
\usepackage{xeCJK}
\usepackage[T1]{fontenc}
\usepackage{xcolor}
\usepackage{setspace}
\setmainfont{Charis SIL}
\setCJKmainfont{Noto Serif CJK HK}
\usepackage{polyglossia}
\newfontfamily\thaifont[Script=Thai]{Noto Serif Thai}
\setmainlanguage{english}
\setotherlanguages{thai}
%\XeTeXlinebreaklocale "th"
%\XeTeXlinebreakskip = 0pt plus 1pt

\title[Pearl River Sinitic]{Sinitic Languages of the Pearl River Region \\ --- Unit 1.3}
\subtitle{13ᵗʰ EACL Summer School, Universität Stuttgart}
\author{Hilário de Sousa 蘇沙}
\institute{EHESS–CRLAO \\ {\tiny(membre associé)}}
\date{Monday 27 July 2026 \\ 丙午六月十四 \\ {\footnotesize \texttt{hilario@bambooradical.com}}}

\begin{document}

\begin{frame}
  \titlepage
\end{frame}

\section{Cantonese Phonotactics}

\begin{frame}[fragile]{Phonotactics}

框 \enquote*{frame}\\
\item {}[kʰ u̯ a ŋ ˥] (TW Mandarin; [˥˩] in BJ Mandarin)\\
\item {}[kʰ u̯ a ŋ ˥] (Cantonese)\\
\bigskip
Initial–Final model:\\ 
{\scriptsize the traditional model used for Mandarin and most Sinitic languages}\\
\medskip
\centering
\begin{tabular}{c|c c c}
    \hline
    Initial & \multicolumn{3}{c}{Final}\\
    & {\tiny medial glide} & \multicolumn{2}{c}{\scriptsize rime}\\
    & & {\scriptsize nucleus} & {\scriptsize coda}\\ \hline
    C & G & V & X\\
    kʰ & u̯ & a & ŋ\\ \hline
\end{tabular}

\end{frame}

\begin{frame}[fragile]{Phonotactics}
e.g. Mandarin Pīnyīn
\begin{columns}
    \begin{column}{0.3\textwidth}
        \begin{figure}
            \includegraphics[width=1\textwidth]{pinyininitial}
        \end{figure}
    \end{column}
    \begin{column}{0.7\textwidth}
        \begin{figure}
            \includegraphics[width=1\textwidth]{pinyinfinal}
        \end{figure}
    \end{column}
\end{columns}

\end{frame}

\begin{frame}[fragile]{Phonotactics}

框 \enquote*{frame}\\
\item {}[kʰ u̯ a ŋ ˥] (TW Mandarin; [˥˩] in BJ Mandarin)\\
\item {}[kʰ u̯ a ŋ ˥] (Cantonese)\\
\bigskip
{\scriptsize often applied to many Yuè and S Pínghuà dialects}\\
Onset–Rime model:\\ 
{\scriptsize (a model more familiar to SE Asia / the West)}\\
\medskip
\centering
\begin{tabular}{c c|c c}
    \hline
    \multicolumn{2}{c|}{Onset} & \multicolumn{2}{c}{Rime}\\
    & & {\scriptsize nucleus} & {\scriptsize coda}\\ \hline
    C & G & V & X\\
    kʰ & u̯ & a & ŋ\\ \hline
\end{tabular}

\end{frame}

\begin{frame}[fragile]{Phonotactics}

Onset–Rime model:\\ 
\medskip
{\centering
\begin{tabular}{c c|c c}
    \hline
    \multicolumn{2}{c|}{Onset} & \multicolumn{2}{c}{Rime}\\
    & & {\scriptsize nucleus} & {\scriptsize coda}\\ \hline
    C & G & V & X\\
    kʰ & u̯ & a & ŋ\\ \hline
\end{tabular} \\}
\medskip
Most Yuè and S Pínghuà dialects are poor in medial glides\\
Standard Cantonese onsets with glides:
\item{} [ku̯- kʰu̯- u̯- i̯-]\\
\medskip
{\scriptsize (Instead of (C)(G)V(X)\textsuperscript{T})} \\
Standard Cantonese usually analysed as \textbf{(C)V(X)}\textsuperscript{T}\\
/kʷ- kʷʰ- w- j-/\\ 
\medskip
\begin{flushright}
\emph{\scriptsize Many people do not distinguish \enquote{Initial–Final} and \enquote{Onset–Rime} \\ both are 聲母–韻母 in Chinese}
\end{flushright}

\end{frame}

\begin{frame}[fragile]{Phonotactics}

The changes from C(G)GVX to \textbf{CVX}:\\
Merger of GV, e.g.\\
\quad 香 *xjaŋ\textsuperscript{A} > Canto. (*hiɔŋ >) /hœŋ˥/\\ 
\quad \quad (cf. HK Hakka /hioŋ˩˧/)\\
\smallskip
Deletion of G, e.g.\\
\quad 菊 *kjuʷk > Canto. /kuk˥/\\ 
\quad \quad (cf. HK Hakka /kʰiuk˧/)\\
\smallskip
Deletion of V, e.g.\\
\quad 偏 *pʰjen\textsuperscript{A} > Canto. /pʰin˥/\\ 
\quad \quad (cf. HK Hakka /pʰien˩˧/)\\
\quad 勸 *kʰwjən\textsuperscript{C} > Canto. /hyn˧/\\
\quad \quad (cf. Nanning Ping. /kʷʰin˨˥/ {\scriptsize V also deleted})\\
\quad \quad (cf. Viet. /xwiə̯n˦\~{}˦˥/ {\scriptsize but /i/ is nucleus})\\
\end{frame}

\section{Romanisation}

\begin{frame}[fragile]{Romanisation}
Cantonese transcription rarely taught in the education system\\
\begin{figure}
    \includegraphics[width=0.85\textwidth]{oldhktextbook}
\end{figure}

\vfill {\tiny \texttt{https://www.hk01.com/社區專題/344752/教科書回憶-收藏殖民地時代課本-消失的科目建構港人中環價值}}\\
\end{frame}

\begin{frame}[fragile]{Romanisation}

There exist many Cantonese Romanisation schemes \\ 
{\scriptsize \emph{(systematic / semi-systematic / ad hoc spelling)}}\\ 
\medskip
Most commonly used nowadays amongst linguists:\\
\textbf{Jyut6 Ping3 粵拼}\\
{\scriptsize (Linguistic Society of Hong Kong Cantonese Romanization Scheme \\ 香港語言學學會粵語拼音方案) \quad \url{https://lshk.org/jyutping-scheme/}}

\bigskip
{\footnotesize See, e.g., \\ 
{\url{https://en.wikipedia.org/wiki/Comparison_of_Cantonese_transcription_systems}} \\ for comparison of Cantonese transcription systems}

\end{frame}

\section{Onsets}

\begin{frame}[fragile]{Onsets}

\begin{columns}
    \begin{column}{0.5\textwidth}
        Onsets in Standard Cantonese\\
        \bigskip
        \begin{tabular}{c c c c c}
            \hline
            \textbf{b} & \textbf{p} & \textbf{m} & \textbf{f} & \\ 
            /p/ & /pʰ/ & /m/ & /f/ & \\ \hline
            \textbf{d} & \textbf{t} & \textbf{n} & & \textbf{l} \\
            /t/ & /tʰ/ & /n/ & & /l/ \\ \hline
            \textbf{z} & \textbf{c} & & \textbf{s} & \textbf{j} \\
            /t͡s/ & /t͡sʰ/ & & /s/ & /j/ \\ \hline
            \textbf{g} & \textbf{k} & \textbf{ng} & \textbf{h} & \\
            /k/ & /kʰ/ & /ŋ/ & /h/ & \\ \hline
            \textbf{gw} & \textbf{kw} & & & \textbf{w} \\
            /kʷ/ & /kʷʰ/ & & & /w/ \\ \hline
        \end{tabular}
    \end{column}
    \begin{column}{0.5\textwidth} 
        \begin{itemize}
            \item zero onset is represented by zero, e.g. 屋 \emph{uk1} \enquote*{house}
            \item \emph{b- p-}, \emph{d- t-} etc. are voiceless (unaspirated vs. aspirated)
            \item some speakers have merged \emph{n-} into \emph{l-}, e.g. 你 \emph{nei5} / \emph{lei5} \enquote*{you (sg)} (vs. 李 \emph{lei5} surname)
            \item \emph{\textbf{j-}} is /j-/ (unusual for Cantonese Romanisations)
        \end{itemize}
    \end{column}
\end{columns}

\end{frame}

\begin{frame}[fragile]{Onsets}

\begin{columns}
    \begin{column}{0.5\textwidth}
        Onsets in Standard Cantonese\\
        \bigskip
        \begin{tabular}{c c c c c}
            \hline
            \textbf{b} & \textbf{p} & \textbf{m} & \textbf{f} & \\ 
            /p/ & /pʰ/ & /m/ & /f/ & \\ \hline
            \textbf{d} & \textbf{t} & \textbf{n} & & \textbf{l} \\
            /t/ & /tʰ/ & /n/ & & /l/ \\ \hline
            \textbf{z} & \textbf{c} & & \textbf{s} & \textbf{j} \\
            /t͡s/ & /t͡sʰ/ & & /s/ & /j/ \\ \hline
            \textbf{g} & \textbf{k} & \textbf{ng} & \textbf{h} & \\
            /k/ & /kʰ/ & /ŋ/ & /h/ & \\ \hline
            \textbf{gw} & \textbf{kw} & & & \textbf{w} \\
            /kʷ/ & /kʷʰ/ & & & /w/ \\ \hline
        \end{tabular}
    \end{column}
    \begin{column}{0.5\textwidth} 
        \begin{itemize}
            \item \emph{z- c- s-} = /t͡s- t͡sʰ s-/
            \item{}[t͡s- t͡sʰ- s-] [t͡ʃ- t͡ʃʰ- ʃ-] not contrastive, with alveolar dominant nowadays 
            \item historically contrastive, cf. 牛池灣 \emph{Ngau \underline{Chi} Wan}, 慈雲山 \emph{\underline{Tsz} Wan \underline{Shan}}, 散石灣 \emph{\underline{San} Shek Wan}
        \end{itemize}
    \end{column}
\end{columns}

\end{frame}

\begin{frame}[fragile]{Onsets}

\begin{columns}
    \begin{column}{0.5\textwidth}
        Onsets in Standard Cantonese\\
        \bigskip
        \begin{tabular}{c c c c c}
            \hline
            \textbf{b} & \textbf{p} & \textbf{m} & \textbf{f} & \\ 
            /p/ & /pʰ/ & /m/ & /f/ & \\ \hline
            \textbf{d} & \textbf{t} & \textbf{n} & & \textbf{l} \\
            /t/ & /tʰ/ & /n/ & & /l/ \\ \hline
            \textbf{z} & \textbf{c} & & \textbf{s} & \textbf{j} \\
            /t͡s/ & /t͡sʰ/ & & /s/ & /j/ \\ \hline
            \textbf{g} & \textbf{k} & \textbf{ng} & \textbf{h} & \\
            /k/ & /kʰ/ & /ŋ/ & /h/ & \\ \hline
            \textbf{gw} & \textbf{kw} & & & \textbf{w} \\
            /kʷ/ & /kʷʰ/ & & & /w/ \\ \hline
        \end{tabular}
    \end{column}
    \begin{column}{0.5\textwidth} 
        \begin{itemize}
            \item older speakers often add \emph{ng-} to 0- (e.g. 亞 \emph{aa3} > \emph{ngaa3}), younger speakers tend to delete \emph{ng-} (e.g. 我 \emph{ngo5} > \emph{o5})
            \item 0- and \emph{ng-} are nearly in complementary distribution: 0- with high tones (1 2 3; e.g. 矮 \emph{ai2} \enquote*{short stature}); \emph{ng-} with low tones (4 5 6; e.g. 蟻 \emph{ngai5} \enquote*{ant})
        \end{itemize}
    \end{column}
\end{columns}

\end{frame}

\begin{frame}[fragile]{Onsets}

\begin{columns}
    \begin{column}{0.5\textwidth}
        Onsets in Standard Cantonese\\
        \bigskip
        \begin{tabular}{c c c c c}
            \hline
            \textbf{b} & \textbf{p} & \textbf{m} & \textbf{f} & \\ 
            /p/ & /pʰ/ & /m/ & /f/ & \\ \hline
            \textbf{d} & \textbf{t} & \textbf{n} & & \textbf{l} \\
            /t/ & /tʰ/ & /n/ & & /l/ \\ \hline
            \textbf{z} & \textbf{c} & & \textbf{s} & \textbf{j} \\
            /t͡s/ & /t͡sʰ/ & & /s/ & /j/ \\ \hline
            \textbf{g} & \textbf{k} & \textbf{ng} & \textbf{h} & \\
            /k/ & /kʰ/ & /ŋ/ & /h/ & \\ \hline
            \textbf{gw} & \textbf{kw} & & & \textbf{w} \\
            /kʷ/ & /kʷʰ/ & & & /w/ \\ \hline
        \end{tabular}
    \end{column}
    \begin{column}{0.5\textwidth} 
        \begin{itemize}
            \item exception 1: 勾鈎鉤 \emph{*gau1} > \emph{ngau1} (taboo avoidance; 勾鈎鉤 \emph{ngau1} contrasting with 歐區 \emph{au1}) 
            \item exception 2: non-sinitic words like 啱 \emph{ngaam1} \enquote*{correct, just, exactly} (cf. Zhuang \emph{ngamq} \enquote*{exactly}), Thai \begin{thai}งาม\end{thai} /ŋaːm˧/ \enquote*{pretty}, Viet \emph{ngám} \enquote*{to fit}) 
        \end{itemize}
    \end{column}
\end{columns}

\end{frame}

\begin{frame}[fragile]{Onsets}

\begin{columns}
    \begin{column}{0.5\textwidth}
        Onsets in Standard Cantonese\\
        \bigskip
        \begin{tabular}{c c c c c}
            \hline
            \textbf{b} & \textbf{p} & \textbf{m} & \textbf{f} & \\ 
            /p/ & /pʰ/ & /m/ & /f/ & \\ \hline
            \textbf{d} & \textbf{t} & \textbf{n} & & \textbf{l} \\
            /t/ & /tʰ/ & /n/ & & /l/ \\ \hline
            \textbf{z} & \textbf{c} & & \textbf{s} & \textbf{j} \\
            /t͡s/ & /t͡sʰ/ & & /s/ & /j/ \\ \hline
            \textbf{g} & \textbf{k} & \textbf{ng} & \textbf{h} & \\
            /k/ & /kʰ/ & /ŋ/ & /h/ & \\ \hline
            \textbf{gw} & \textbf{kw} & & & \textbf{w} \\
            /kʷ/ & /kʷʰ/ & & & /w/ \\ \hline
        \end{tabular}
    \end{column}
    \begin{column}{0.5\textwidth} 
        \begin{itemize}
            \item \emph{h-} is glottal [h] (not velar [x] like Mandarin)
            \item some people merge \emph{gwo kwo} into \emph{go ko} e.g. 國 \emph{gwok3} > \emph{gok3} (= 各 \emph{gok3})
        \end{itemize}
    \end{column}
\end{columns}

\end{frame}

\section{Rimes}

\begin{frame}[fragile]{Codas}

\textbf{-p -m -t -n -k -ng -i -u}\\
/-p -m -t -n -k -ŋ -j -w/\\

\bigskip
\begin{itemize}
    \item /-p -t -k/: unreleased [p̚  t̚  k̚]
    \item highly stigmatised: errors with /-p -m/ 
    \item stigmatised: total or partial merger of /-k -ŋ/ into /-t -n/\\
    e.g. 香港 [hœŋ˥ kɔŋ˨˥] > [hœn˥ kɔn˨˥], 北角 [pɐk˥ kɔk˧] > [pɐt˥ kɔt˧]
    \item newer minor trends: \enquote{indeterminate} merger of /-n -ŋ/ as something like Ṽ, and /-t -k/ as something like ʔ\\
    e.g. 公安 [kʊŋ˥ ɔn˥] > [kʊ̃ ɔ̃], 失約 [sɐt˥ jœk˧] > [sɐʔ jœʔ]
\end{itemize}

\end{frame}

\begin{frame}[fragile]{Codas}

\textbf{-p -m -t -n -k -ng -i -u}\\
/-p -m -t -n -k -ŋ -j -w/\\

\bigskip
(Stereotypical) non-native accents, e.g.:
\begin{itemize}
    \item /-m/ > /-n/: e.g. Mandarin, Wu
    \item /-t -n/ > /-k -ŋ/: e.g. 潮州 Teochew
    \item /-p -m -t -n/ {\tiny partial or total} > /-k -ŋ/: \\莞寶粵 e.g. Gun2Bou2 Jyut6 / GuǎnBǎo Yuè 
    
\end{itemize}

\end{frame}

\begin{frame}[fragile]{Vowels}

\begin{figure}
    \includegraphics[width=0.85\textwidth]{cantovowel}
\end{figure}

\vfill {\tiny \setstretch{0.1} Based on Zee, Eric. 1999. Chinese (Hong Kong Cantonese). In The International Phonetic Association (ed.), \emph{Handbook of the International Phonetic Association}, 58–60. Cambridge: Cambridge University Press.\\ Modified from \url{https://upload.wikimedia.org/wikipedia/commons/c/c6/Cantonese_vowel_chart.svg}\par}

\end{frame}

\begin{frame}[fragile]{Rimes}

\begin{table}
\scriptsize
\begin{tabular}{c|c c c c c c c c c}
    \hline
     & -0 & /-p/ & /-m/ & /-t/ & /-n/ & /-k/ & /-ŋ/ & /-j/ & /-w/\\ \hline
    /i/ & \textbf{i} & \textbf{ip} & \textbf{im} & \textbf{it} & \textbf{in} & \textbf{ik} & \textbf{ing} & &\textbf{iu}\\
     & [i] & [ip] & [im] & [it] & [in] & [ek] & [eŋ] & & [iu]\\ \hline
    /u/ & \textbf{u} & & & \textbf{ut} & \textbf{un} & \textbf{uk} & \textbf{ung} & \textbf{ui} & \\
     & [u] & & & [ut] & [un] & [ok] & [oŋ] & [ui] & \\ \hline
    /y/ & \textbf{yu} & & & \textbf{yut} & \textbf{yun} & & & & \\
     & [y] & & & [yt] & [yn] & & & & \\ \hline
    /ɛ/ & \textbf{e} & \textbf{ep} & \textbf{em} & \textbf{(et)} & \textbf{(en)} & \textbf{ek} & \textbf{eng} & \textbf{ei} & \textbf{eu} \\
     & [ɛ] & [ɛp] & [ɛm] & ([ɛt]) & ([ɛn]) & [ɛk] & [ɛŋ] & [ei] & [ɛu] \\ \hline
    /ɔ/ & \textbf{o} & & & \textbf{ot} & \textbf{on} & \textbf{ok} & \textbf{ong} & \textbf{oi} & \textbf{ou} \\
     & [ɔ] & & & [ɔt] & [ɔn] & [ɔk] & [ɔŋ] & [ɔi] & [ou] \\ \hline
    /œ/ & \textbf{oe} & & & \textbf{(oet)} & & \textbf{oek} & \textbf{oeng} & & \\
     & [œ] & & & ([œt]) & & [œk] & [œŋ] & & \\
     & & & & \textbf{eot} & \textbf{eon} & & & \textbf{eoi} & \\
     & & & & [ɵt] & [ɵn] & & & [ɵy] & \\ \hline
    /ɐ/ & \textbf{(a)} & \textbf{ap} & \textbf{am} & \textbf{at} & \textbf{an} & \textbf{ak} & \textbf{ang} & \textbf{ai} & \textbf{au} \\
     & ([ɐ]) & [ɐp] & [ɐm] & [ɐt] & [ɐn] & [ɐk] & [ɐŋ] & [ɐi] & [ɐu] \\ \hline
    /a/ & \textbf{aa} & \textbf{aap} & \textbf{aam} & \textbf{aat} & \textbf{aan} & \textbf{aak} & \textbf{aang} & \textbf{aai} & \textbf{aau} \\
     & [a] & [ap] & [am] & [at] & [an] & [ak] & [aŋ] & [ai] & [au] \\ \hline
    0 & & & \textbf{m} & & & & \textbf{ng} & & \\
     & & & [m̩] & & & & [ŋ̍] & & \\ \hline
\end{tabular}
\end{table}

\end{frame}

\end{document}
